\documentclass[a4paper,slidestop,xcolor=pst,blue]{beamer}

\usepackage{beamerthemesplit}
\usepackage[utf8]{inputenc}
\usepackage[spanish]{babel}
\usepackage{graphicx}
\usepackage{pstricks} % PSTricks package
\usepackage{setspace}
\usepackage{multirow}
\usepackage{listings}
\usepackage{pgfpages}
\usepackage{hyperref}
\usepackage{etoolbox}
\usepackage{epstopdf}

\makeatletter
\patchcmd{\beamer@sectionintoc}{\vskip1.5em}{\vskip0.5em}{}{}
\makeatother

\setbeamercovered{dynamic}
\setcounter{tocdepth}{2}
\setbeamercolor{frametitle}{fg=black,bg=white}
\setbeamercolor{section in toc shaded}{fg=black}
\setbeamercolor{section in toc}{fg=red}
\setbeamercolor{subsection in toc shaded}{fg=black}
\setbeamercolor{subsection in toc}{fg=red}
\setbeamerfont{section in toc}{size=\small}
\setbeamerfont{subsection in toc}{size=\small}
\setbeamertemplate{section in toc shaded}[default][99]
\setbeamertemplate{subsection in toc shaded}[default][99]

\AtBeginSection[]
{\begin{frame}[c]
  \frametitle{Índice}
	\tableofcontents[currentsection,
        sectionstyle=show/shaded,
        subsectionstyle=hide]
\end{frame}}

\AtBeginSubsection[]
{\begin{frame}[c]
	\frametitle{Índice}
	\tableofcontents[
  		currentsection,
  		sectionstyle=shaded/shaded,
  		currentsubsection,
  		subsectionstyle=show/shaded/hide
		]
\end{frame}}

\setbeamercolor{frametitle}{fg=black,bg=white}

\setbeamertemplate{frametitle}{
	\begin{centering}
		\insertframetitle
		\par
	\end{centering}
}

\usetheme[secheader]{Boadilla} 

\title[Procesos de Negocio]{Especificación y Modelado de Procesos de Negocio}

\author[P. S{\'a}nchez]{\alert{Pablo S{\'a}nchez}}

\institute[IIE]{
		   Dpto. Ingenier{\'i}a Inform{\'a}tica y Electr{\'o}nica \\
		   Universidad de Cantabria \\
		   Santander (Cantabria, Espa{\~n}a) \\
		   \texttt{p.sanchez@unican.es}
}

\date{}

\begin{document}

\begin{frame}[c]
	\titlepage
	\begin{columns}
		\column{0.50\linewidth}
			\centering
    		\includegraphics[width=.28\textwidth,keepaspectratio=true]{images/istr.eps}
		\column{0.50\linewidth}
			\centering
			\includegraphics[width=.25\textwidth,keepaspectratio=true]{images/uc.eps}
	\end{columns}
\end{frame}

\begin{frame}[c]
    \frametitle{\alert{Advertencia}}
    \begin{center}
        Todo el material contenido en este documento no constituye en modo alguno una obra de referencia o apuntes oficiales mediante el cual se puedan preparar las pruebas evaluables necesarias para superar la asignatura. \ \\
        \ \\
        Este documento contiene exclusivamente una serie de diapositivas cuyo objetivo es servir de complemento visual a las actividades realizadas en el aula para la transmisi{\'o}n del contenido sobre el cual versar{\'a}n las mencionadas pruebas evaluables.  \ \\
        \ \\
        Dicho de forma m{\'a}s clara, \alert{estas transparencias no son apuntes y su objetivo no es servir para que el alumno pueda preparar la asignatura.}
    \end{center}
\end{frame}

\section{Introducción}

\subsection{Objetivos y Bibliografía}

\begin{frame}[c]
    \frametitle{Objetivos del Tema}
    \begin{enumerate}[<+->]
         \item Comprender qué es un proceso de negocio.
         \item Conocer qué es BPMN 2.0.
         \item Comprender cómo se utilizan los principales elementos de BPMN 2.0.
         \item Ser capaz de modelar procesos de negocio básicos en BPMN 2.0.
    \end{enumerate}
\end{frame}

\begin{frame}[c]
    \frametitle{Bibliografía}
    \begin{thebibliography}{1}
        
        %% Estander
        
        %% Buscar libro en biblioteca
    
    \end{thebibliography}
\end{frame}

\section{Proceso de Negocio}

%% Definición

\section{BPMN 2.0}

%% Definición y características.

\begin{frame}[c]
    \frametitle{Elementos BPMN 2.0}
    %% Participants, Pools y Lanes
    %% Tasks
    %% Start Event, End Event?
    %% Flow
    %% Parallel Gateway, Exclusive Gateway, Inclusive Gateway, Synchronizing Gateway
\end{frame}



\section{Sumario y Conclusiones}

\begin{frame}[c]
    \frametitle{Sumario}
    %% Definición
\end{frame}

\begin{frame}[c]
    \frametitle{Conclusiones}
    %% Definición
\end{frame}


\end{document}
